\documentclass[11pt,a4paper]{article}
\usepackage[T1]{fontenc}
\usepackage[utf8]{inputenc}
\usepackage{lmodern}
\usepackage{amsmath,amssymb,amsthm,mathtools,mathrsfs}
\usepackage{graphicx,float,xcolor,multicol}
\usepackage[shortlabels]{enumitem}
\usepackage[margin=27mm]{geometry}
\usepackage{microtype,needspace,placeins}
\usepackage{tikz,tikz-cd}
\usetikzlibrary{arrows.meta,calc,positioning,patterns,decorations.pathreplacing}
\usepackage[colorlinks=true,linkcolor=teal,urlcolor=teal,citecolor=teal]{hyperref}
\setlength{\parindent}{0pt}
\setlength{\parskip}{.6em}
\setcounter{secnumdepth}{0}
\allowdisplaybreaks
\raggedbottom
\providecommand{\cross}{\times}
\DeclareUnicodeCharacter{2020}{\ensuremath{\dagger}}
\DeclareMathOperator{\supp}{supp}
\DeclareMathOperator*{\esssup}{ess\,sup}
\providecommand{\chapter}{\section}
\newenvironment{tool}[2]{\subsection*{Tool #1 --- #2}}{}
\newcommand{\ToolRef}[1]{Tool #1}
\newcommand{\FigureTag}[2]{\textit{Figure #1}}
\newcommand{\Result}[1]{\subsection*{#1}}
\newcommand{\Application}[1]{\subsubsection*{Application\if\relax\detokenize{#1}\relax\else\space--- #1\fi}}
\newcommand{\ProofHeading}{\par\textit{Proof.}\space}
\newcommand{\ProofOf}[1]{\par\textit{Proof of #1.}\space}
\newcommand{\RemarkHeading}{\par\textbf{Remark.}\space}
\newcommand{\NoteHeading}{\par\textbf{Note.}\space}
\newcommand{\ExerciseHeading}{\par\textbf{Exercise.}\space}

% Rendering support only; mathematical text is in each independent post.
\usepackage{booktabs,array,longtable}
\definecolor{HandoutBlue}{HTML}{2F6690}
\definecolor{HandoutNavy}{HTML}{17365D}
\newtheorem{theorem}{Theorem}
\newtheorem{lemma}[theorem]{Lemma}
\newtheorem{proposition}[theorem]{Proposition}
\newtheorem{corollary}[theorem]{Corollary}
\newtheorem{conjecture}[theorem]{Conjecture}
\newtheorem{definition}{Definition}
\newtheorem{example}{Example}
\newtheorem{namedtoolcounter}{Tool}
\newenvironment{namedtool}[1]{\begin{namedtoolcounter}[#1]}{\end{namedtoolcounter}}
\newtheorem*{claim}{Claim}
\newtheorem*{remark}{Remark}
\newtheorem*{exercise}{Exercise}
\newtheorem*{question}{Question}
\newtheorem*{observation}{Observation}
\newcommand{\SourceChapter}[2]{\section*{#2}}
\newcommand{\Lecture}[2]{\section*{Lecture #1: #2}}
\newcommand{\unclear}[1]{\textit{[unclear: #1]}}
\newcommand{\sourcegap}{\textit{[unfinished in the source]}}
\DeclareMathOperator{\dist}{dist}
\DeclareMathOperator{\id}{id}
\DeclareMathOperator{\Hol}{Hol}
\DeclareMathOperator{\Per}{Per}
\DeclareMathOperator{\Fix}{Fix}
\DeclareMathOperator{\diam}{diam}
\DeclareMathOperator{\Var}{Var}
\DeclareMathOperator{\Leb}{Leb}
\DeclareMathOperator{\Hom}{Hom}
\DeclareMathOperator{\End}{End}
\DeclareMathOperator{\Ann}{Ann}
\DeclareMathOperator{\Spec}{Spec}
\DeclareMathOperator{\Max}{Max}
\DeclareMathOperator{\Ob}{Ob}
\DeclareMathOperator{\Tor}{Tor}
\DeclareMathOperator{\Ass}{Ass}
\DeclareMathOperator{\Mor}{Mor}
\DeclareMathOperator{\Fun}{Fun}
\DeclareMathOperator{\Nat}{Nat}
\DeclareMathOperator{\Cone}{Cone}
\DeclareMathOperator{\MCG}{MCG}
\DeclareMathOperator{\Aut}{Aut}
\DeclareMathOperator{\Deck}{Deck}
\DeclareMathOperator{\SL}{SL}
\DeclareMathOperator{\PSL}{PSL}
\DeclareMathOperator{\Area}{Area}
\DeclareMathOperator{\im}{im}
\DeclareMathOperator{\PSh}{PSh}
\newcommand{\CRing}{\mathsf{CRings}}
\newcommand{\AMod}{A\text{-}\mathsf{Mod}}
\newcommand{\Ab}{\mathsf{Ab}}
\newcommand{\Z}{\mathbb Z}
\newcommand{\Q}{\mathbb Q}
\newcommand{\R}{\mathbb R}
\newcommand{\C}{\mathbb C}
\newcommand{\N}{\mathbb N}
\newcommand{\kk}{\mathbb k}
\renewcommand{\aa}{\mathfrak a}
\newcommand{\bb}{\mathfrak b}
\newcommand{\pp}{\mathfrak p}
\newcommand{\qq}{\mathfrak q}
\newcommand{\mm}{\mathfrak m}
\newcommand{\nn}{\mathfrak n}
\newcommand{\rad}{\mathrm r}
\newcommand{\cat}[1]{\mathsf{#1}}
\setcounter{secnumdepth}{0}
\allowdisplaybreaks

\graphicspath{{figures/ergodic-theory/}}
\begin{document}
\section*{Mean, Maximal, and Pointwise Ergodic Theorems}
% BLOG-CONTENT-BEGIN
In our attempt to get a hold on the nature of $\{T^{\circ n}(x)\}_n$ or
$\{T^{\circ n}(E)\}_n$, we saw that $\{T^{\circ n}(x)\}_n$ is recurrent.
Further, if $T$ is ergodic, then $\{T^{\circ n}(x)\}_n$ is recurrent on all
positive-measure subsets of $X$ (for almost every $x\in X$). The next natural
question to ask is: how recurrent are they? Mathematically speaking, we ask
what happens, as $N\to\infty$, to
\[
  A_N^f:=\frac1N\sum_{n=0}^{N-1}U_T^n f
  \qquad\text{(the $T$-average of $f$).}
\]
In answering what happens to $\frac1{N+1}\sum_{n=0}^{N}U_T^n f$ as
$N\to\infty$, ergodic theorems will, in some sense, establish a duality
between space and time.

\setcounter{theorem}{8}\begin{theorem}[Birkhoff]
Let $(X,\mathcal B,\mu,T)$ be a measure-preserving system, with $T$ ergodic.
If $f\in L^1_\mu$, then
\[
  \lim_{n\to\infty}\frac1n\sum_{j=0}^{n-1}U_T^j f=\int f\,d\mu.
\]
\end{theorem}

% Source page 2
The left-hand side is an average over time, and the right-hand side, being
an integral, is an average over space. If $f=\chi_B$, then the ergodic
theorem is telling you that the probability of $T^n(x)\in B$ is given by
the size of $B$ (or the measure of $B$). The independence of the theorem
from $x\in X$ is the remarkable thing to observe, and the imposition that
$T$ be ergodic is essential; otherwise, there may be $x\in X$ such that
$T^n(x)\notin B$ for all $n\geq1$. But as soon as you remove that
set-theoretic hindrance, recurrence of orbits is governed by the sizes of
sets.


Thus certain spatial information can be translated into temporal information,
and vice versa.

% Source page 3

\setcounter{theorem}{9}\begin{theorem}[von Neumann's mean ergodic theorem: $L^2_\mu$ version]
\label{thm:l3-mean-l2}
Let $(X,\mathcal B,\mu,T)$ be a measure-preserving system. Let $P_T$ be the
orthogonal projection of $L^2_\mu$ onto
\[
  I=\{g\in L^2_\mu:U_Tg=g\}.
\]
Then, for every $f\in L^2_\mu$,
\[
  \frac1N\sum_{n=0}^{N-1}U_T^nf\longrightarrow P_Tf
  \quad\text{in the $L^2_\mu$ norm.}
\]
\end{theorem}

\setcounter{theorem}{10}\begin{theorem}[von Neumann's mean ergodic theorem: $L^1_\mu$ version]
\label{thm:l3-mean-l1}
Let $(X,\mathcal B,\mu,T)$ be a measure-preserving system. For every
$f\in L^1_\mu$,
\[
 \frac1N\sum_{n=0}^{N-1}U_T^nf\longrightarrow f'\in L^1_\mu
 \quad\text{in $L^1_\mu$,}\qquad U_Tf'=f'.
\]
\end{theorem}

% Source page 6
\begin{remark}
\label{rem:l3-frequency}
If $T$ is ergodic, then, for $f=\chi_B$, we only know so far that the
frequency of orbit visits of $f$ is a fixed number, independent of $x$.
\end{remark}

\setcounter{theorem}{11}\begin{theorem}[Maximal ergodic theorem]
\label{thm:l3-maximal}
Let $(X,\mathcal B,\mu,T)$ be a measure-preserving system, and let
$g:X\to\mathbb R$, $g\in L^1_\mu$. Define, for $\alpha\in\mathbb R$,
\[
 E_\alpha:=\left\{x\in X:
 \sup_{n\geq1}\frac1n\sum_{i=0}^{n-1}g(T^ix)>\alpha\right\}.
\]
Then
\[
 \alpha\mu(E_\alpha)\leq\int_{E_\alpha}g\,d\mu\leq\|g\|_1,
 \qquad
 \alpha\mu(E_\alpha\cap A)\leq\int_{E_\alpha\cap A}g\,d\mu
 \quad\text{if }T^{-1}(A)=A.
\]
\end{theorem}

% Source page 7
We shall first prove a technical lemma and derive the theorem from it in a snap.
\setcounter{theorem}{12}\begin{lemma}
\label{lem:l3-positive}
Let $U:L^1_\mu\to L^1_\mu$ be a positive linear operator with
$\|U\|\leq1$. Given $f\in L^1_\mu$, set
\[
 f_k=\sum_{\ell=0}^{k-1}U^\ell f,\qquad f_0=0,
 \qquad F_N=\max\{f_n:0\leq n\leq N\}.
\]
Then $\int_{\{x:F_N(x)>0\}}f\,d\mu\geq0$ for every $N\geq1$.
\end{lemma}
\begin{proof}
We have $F_N\in L^1_\mu$, since $F_N\leq\sum_{n=0}^N|f_n|$.
Also, $F_N\geq f_n$ for $0\leq n\leq N$, so positivity and linearity give
$UF_N+f\geq Uf_n+f=f_{n+1}$. Thus
$UF_N+f\geq\max_{1\leq n\leq N}f_n$.
Set $P=\{x:F_N(x)>0\}$. For $x\in P$,
\[
 F_N(x)=\max_{0\leq n\leq N}f_n(x)
       =\max_{1\leq n\leq N}f_n(x),
\]
since $f_0=0$. Hence $UF_N(x)+f(x)\geq F_N(x)$, or
$f(x)\geq F_N(x)-UF_N(x)$. By positivity, $UF_N(x)\geq0$, as
$F_N(x)\geq0$. It follows that
\[
 \int_P f\,d\mu
 \geq\int_P F_N\,d\mu-\int_P UF_N\,d\mu
 =\int_XF_N\,d\mu-\int_PUF_N\,d\mu,
\]
% Source page 8
because $F_N(x)=0$ if $x\notin P$. Therefore
\[
 \int_Pf\,d\mu\geq\int_XF_N\,d\mu-\int_XUF_N\,d\mu
 =\|F_N\|_1-\|UF_N\|_1\geq0.
\]
Here the equality uses positivity, and the last inequality uses $\|U\|\leq1$.
\end{proof}
\begin{proof}[Proof of Theorem~\href{https://codezen1729.github.io/math-blog/\#/post/mean-maximal-and-pointwise-ergodic-theorems}{12}]
Set $f=g-\alpha$ and $Uf:=f\circ T$. This is a unitary positive operator,
so $\|U\|\leq1$. We have
\[
 E_\alpha=\bigcup_{N=0}^\infty\{x:F_N(x)>0\}.
\]
By the lemma,
$\int_{E_\alpha}f\,d\mu\geq0$, so
$\int_{E_\alpha}g\,d\mu\geq\alpha\mu(E_\alpha)$.
To obtain the other inequality, look at
\[
 \bigl(A,\mathcal B|_A,\mu(A)^{-1}\mu|_A,T|_A\bigr).\qedhere
\]
\end{proof}


\setcounter{theorem}{13}\begin{theorem}[Birkhoff's ergodic theorem]
\label{thm:l3-birkhoff}
If $f\in L^1_\mu$, then
\[
 \lim_{n\to\infty}\frac1n\sum_{j=0}^{n-1}f(T^j(x))=f^*(x).
\]
The function $f^*$ is $T$-invariant, and $\int f^*\,d\mu=\int f\,d\mu$.
\end{theorem}

% Source page 9

\begin{remark}[Final remark]
If $T$ is ergodic, then $f^*$ must be constant, since it is $T$-invariant
(the eigenfunctions of $U_T$ corresponding to $1$ are only the constants).
Thus
\[
 \int f^*\,d\mu=f^*\mu(X)=f^*=\int f\,d\mu,
\]
and we finally get the link between space and time averages:
\[
 \boxed{\lim_{n\to\infty}\frac1n\sum_{i=0}^{n-1}f(T^i(x))
 =\int_Xf\,d\mu.}
\]
\end{remark}
% BLOG-CONTENT-END
\end{document}
